\chapter{Lattice Yang--Mills as a Nonperturbative Formulation}
\label{ch:lattice-yang-mills-nonperturbative-definition}
\ConventionsEuclidean

This chapter gives a regulator-level construction of pure Yang--Mills theory
which does not begin by choosing a covariant gauge.  The construction is a
finite-dimensional integral over compact groups at finite lattice spacing and
finite volume.  Its continuum limit is a separate problem: one must prove, or
assume with stated evidence, that a sequence of finite-cutoff theories has
gauge-invariant Schwinger functions with a limit satisfying the Euclidean
axioms and yielding a positive-energy Hilbert-space theory after
Osterwalder--Schrader reconstruction.

\section{Finite Lattice Gauge Fields}

Fix a compact Lie group \(G\).  For most formulas below \(G=SU(N)\), and
\(\operatorname{tr}\) denotes the trace in the fundamental representation with
\(\operatorname{tr}(t^a t^b)=\delta^{ab}\), matching the continuum
Yang--Mills convention of
Chapter~\ref{ch:vol2-classical-yang-mills-matter}.  Let
\[
  \Lambda_{a,L}
  =
  (a\mathbb Z/L\mathbb Z)^4
\]
be a four-dimensional periodic hypercubic lattice, where \(a>0\) is the
lattice spacing and \(L/a\) is an integer.  The positively oriented links are
\[
  E^+(\Lambda_{a,L})
  =
  \{(x,\mu)\mid x\in\Lambda_{a,L},\ \mu=1,\ldots,4\}.
\]
The link \((x,\mu)\) joins \(x\) to \(x+a\widehat\mu\).  A lattice gauge field
is a function
\[
  U:E^+(\Lambda_{a,L})\longrightarrow G,
  \qquad
  (x,\mu)\longmapsto U_\mu(x).
\]
The inverse orientation is assigned the inverse group element,
\[
  U_{-\mu}(x+a\widehat\mu)=U_\mu(x)^{-1}.
\]
Thus a lattice gauge field is a point of the compact manifold
\[
  \mathcal A_{\Lambda}=G^{E^+(\Lambda_{a,L})}.
\]

The lattice gauge group is
\[
  \mathcal G_{\Lambda}=G^{\Lambda_{a,L}}.
\]
An element \(\Omega\in\mathcal G_{\Lambda}\) acts by
\begin{equation}
  U_\mu(x)\longmapsto
  U^\Omega_\mu(x)
  =
  \Omega(x)U_\mu(x)\Omega(x+a\widehat\mu)^{-1}.
  \label{eq:lattice-link-gauge-transformation}
\end{equation}
This convention matches the Hermitian continuum convention of
Chapter~\ref{ch:vol2-classical-yang-mills-matter}.  For a smooth field with
\[
  U_\mu(x)=\exp\bigl(\ii a A_\mu(x)+O(a^2)\bigr),
\]
the endpoint rule \eqref{eq:lattice-link-gauge-transformation} gives
\[
  A_\mu^\Omega
  =
  \Omega A_\mu\Omega^{-1}
  +
  \ii\,\Omega\partial_\mu\Omega^{-1}
  +O(a).
\]

For \(\mu<\nu\), define the oriented plaquette based at \(x\) by
\begin{equation}
  U_{\mu\nu}(x)
  =
  U_\mu(x)
  U_\nu(x+a\widehat\mu)
  U_\mu(x+a\widehat\nu)^{-1}
  U_\nu(x)^{-1}.
  \label{eq:lattice-plaquette-definition}
\end{equation}
Under \eqref{eq:lattice-link-gauge-transformation},
\begin{equation}
  U_{\mu\nu}(x)\longmapsto
  \Omega(x)U_{\mu\nu}(x)\Omega(x)^{-1}.
  \label{eq:lattice-plaquette-covariance}
\end{equation}
Therefore every class function of \(U_{\mu\nu}(x)\), and in particular
\(\operatorname{Re}\operatorname{tr}U_{\mu\nu}(x)\), is gauge invariant.

\begin{figure}[H]
  \centering
  \resizebox{0.82\textwidth}{!}{%
  \begin{tikzpicture}[x=1cm,y=1cm,>=Stealth,every node/.style={font=\scriptsize}]
    \coordinate (x) at (0,0);
    \coordinate (xm) at (3.2,0);
    \coordinate (xn) at (0,2.2);
    \coordinate (xmn) at (3.2,2.2);
    \fill[black] (x) circle (2pt);
    \fill[black] (xm) circle (2pt);
    \fill[black] (xn) circle (2pt);
    \fill[black] (xmn) circle (2pt);
    \draw[->,very thick,blue!70!black] (x)--(xm)
      node[midway,below] {\(U_\mu(x)\)};
    \draw[->,very thick,blue!70!black] (xm)--(xmn)
      node[midway,right] {\(U_\nu(x+a\widehat\mu)\)};
    \draw[->,very thick,blue!70!black] (xmn)--(xn)
      node[midway,above] {\(U_\mu(x+a\widehat\nu)^{-1}\)};
    \draw[->,very thick,blue!70!black] (xn)--(x)
      node[midway,left] {\(U_\nu(x)^{-1}\)};
    \node[below left] at (x) {\(x\)};
    \node[below right] at (xm) {\(x+a\widehat\mu\)};
    \node[above left] at (xn) {\(x+a\widehat\nu\)};
    \node[above right] at (xmn) {\(x+a\widehat\mu+a\widehat\nu\)};
    \node[align=center] at (1.6,1.1)
      {\(U_{\mu\nu}(x)\)\\conjugates at \(x\)};

    \begin{scope}[shift={(5.2,0.25)}]
      \node[draw,rounded corners=2pt,inner sep=5pt,align=center] (links)
        at (0,1.45) {link variables\\\(U_\mu(x)\in G\)};
      \node[draw,rounded corners=2pt,inner sep=5pt,align=center] (plaq)
        at (3.1,1.45) {plaquette\\\(U_{\mu\nu}(x)\)};
      \node[draw,rounded corners=2pt,inner sep=5pt,align=center] (act)
        at (6.0,1.45) {class functions\\\(\operatorname{Re}\operatorname{tr}U_{\mu\nu}\)};
      \draw[->,thick] (links)--(plaq);
      \draw[->,thick] (plaq)--(act);
      \node[align=center] at (3.05,0.05)
        {gauge covariance is exact\\at finite lattice spacing};
    \end{scope}
  \end{tikzpicture}
  }
  \caption{The plaquette is the elementary lattice curvature.  Link variables
  transform at their endpoints; around a closed elementary square the endpoint
  factors cancel except for conjugation at the basepoint.}
  \label{fig:lattice-plaquette-gauge-covariance}
\end{figure}

\section{Wilson Action and the Finite Haar Integral}

For \(G=SU(N)\) the Wilson action is
\begin{equation}
  S_{W,\Lambda}[U]
  =
  \frac{\beta}{N}
  \sum_{x\in\Lambda_{a,L}}
  \sum_{1\leq\mu<\nu\leq 4}
	  \operatorname{Re}\operatorname{tr}
	  \bigl(\mathbf 1-U_{\mu\nu}(x)\bigr),
	  \qquad
	  \beta=\frac{N}{g_0^2}.
	  \label{eq:lattice-wilson-action}
\end{equation}
Here \(g_0\) is the lattice coupling parameter in the trace-form continuum
normalization
\[
  S_{\mathrm{cont}}[A]
  =
  \frac1{4g_0^2}\int\operatorname{tr}(F_{\mu\nu}F_{\mu\nu}) .
\]
With the more common generator normalization
\(\operatorname{tr}(t^a t^b)=\frac12\delta^{ab}\), the same Wilson action is
often written with \(\beta=2N/g_0^2\).  The difference is a coordinate change in
the generator and coupling normalization, not a different lattice integral.
The integrand
\(\operatorname{Re}\operatorname{tr}(\mathbf 1-U_{\mu\nu})\) is nonnegative,
because each eigenvalue of \(U_{\mu\nu}\) lies on the unit circle and hence
\(\operatorname{Re}\operatorname{tr}U_{\mu\nu}\leq N\).  Gauge invariance of
\(S_{W,\Lambda}\) follows from
\eqref{eq:lattice-plaquette-covariance}.

Let \(\dd U\) be normalized Haar measure on \(G\).  At finite \(a\) and finite
\(L\), the lattice partition function is the finite-dimensional compact
integral
\begin{equation}
  Z_{\Lambda}(\beta)
  =
  \int_{\mathcal A_{\Lambda}}
  \prod_{\ell\in E^+(\Lambda_{a,L})}\dd U_\ell\,
  \exp\{-S_{W,\Lambda}[U]\}.
  \label{eq:lattice-partition-function}
\end{equation}
Since \(\mathcal A_{\Lambda}\) is compact and the integrand is continuous,
\(Z_{\Lambda}(\beta)\) is finite and strictly positive.  For a bounded
gauge-invariant function \(O:\mathcal A_{\Lambda}\to\mathbb C\),
\begin{equation}
  \langle O\rangle_{\Lambda,\beta}
  =
  Z_{\Lambda}(\beta)^{-1}
  \int_{\mathcal A_{\Lambda}}
  \prod_{\ell}\dd U_\ell\,
  O[U]\,\exp\{-S_{W,\Lambda}[U]\}
  \label{eq:lattice-expectation}
\end{equation}
is a completely defined number.  No division by an infinite gauge-group volume
and no Faddeev--Popov determinant enter this finite-cutoff definition.

The relation to the continuum Yang--Mills action is obtained by evaluating the
plaquette on a smooth continuum connection.  With
\[
  U_\mu(x)
  =
  \exp\bigl(\ii a A_\mu(x+\tfrac a2\widehat\mu)+O(a^3)\bigr),
\]
the Baker--Campbell--Hausdorff expansion gives
\begin{equation}
  U_{\mu\nu}(x)
  =
  \exp\bigl(\ii a^2 F_{\mu\nu}(x)+O(a^3)\bigr),
  \label{eq:lattice-plaquette-continuum-expansion}
\end{equation}
where \(F_{\mu\nu}\) is the Hermitian-coordinate curvature of
Chapter~\ref{ch:vol2-classical-yang-mills-matter}.  For \(SU(N)\),
\(\operatorname{tr}F_{\mu\nu}=0\), and hence
\[
  \operatorname{Re}\operatorname{tr}(\mathbf 1-U_{\mu\nu})
  =
  \frac{a^4}{2}\operatorname{tr}(F_{\mu\nu}F_{\mu\nu})
  +O(a^6).
\]
Substitution in \eqref{eq:lattice-wilson-action} gives
\[
  \sum_{\mu<\nu}\frac{a^4}{2}\operatorname{tr}(F_{\mu\nu}F_{\mu\nu})
  =
  \frac{a^4}{4}
  \sum_{\mu,\nu}\operatorname{tr}(F_{\mu\nu}F_{\mu\nu}),
\]
because \(F_{\nu\mu}=-F_{\mu\nu}\).  Therefore
\begin{equation}
  S_{W,\Lambda}[U(A)]
  =
  \frac{1}{4g_0^2}
  \int_{\mathbb T_L^4}
  \operatorname{tr}(F_{\mu\nu}F_{\mu\nu})\,\dd^4x
  +O(a^2),
  \label{eq:lattice-wilson-continuum-limit-formal}
\end{equation}
for smooth fields on the torus of side \(L\).  Equation
\eqref{eq:lattice-wilson-continuum-limit-formal} is an expansion of the
classical lattice action on smooth configurations.  The quantum continuum
limit also requires control of the measure as \(a\downarrow0\), \(L\) is sent
to the desired infrared limit, and \(\beta\to\infty\) along a trajectory that
keeps a physical scale fixed.

\section{Reflection Positivity and the Transfer Matrix}

The Wilson lattice theory is Euclidean from the start.  To connect it to a
Hilbert-space theory one uses the lattice version of reflection positivity,
the discrete precursor of the Osterwalder--Schrader condition developed in
Chapter~\ref{ch:os-reconstruction}.

Choose the reflection \(\vartheta\) which sends the Euclidean time coordinate
\(x_4\) to \(-x_4\) and reverses the orientation of reflected links.  Let
\(\mathcal B_+\) be the algebra generated by gauge-invariant functions of link
variables whose support lies in the positive-time half lattice.  For
\(F\in\mathcal B_+\), set
\[
  \Theta F[U]=\overline{F[\vartheta U]}.
\]

\begin{theorem}[Osterwalder--Seiler reflection positivity for the Wilson action]
\label{thm:lattice-wilson-reflection-positivity}
On a reflection-invariant finite hypercubic lattice, the Wilson measure
\[
  \dd\mu_{\Lambda,\beta}(U)
  =
  Z_{\Lambda}(\beta)^{-1}
  \prod_{\ell}\dd U_\ell\,\ee^{-S_{W,\Lambda}[U]}
\]
satisfies
\begin{equation}
  \int \Theta F[U]\,F[U]\,\dd\mu_{\Lambda,\beta}(U)\geq0,
  \qquad
  F\in\mathcal B_+.
  \label{eq:lattice-reflection-positivity}
\end{equation}
\end{theorem}

\begin{proof}
The proof is finite-dimensional.  Haar measure is invariant under reflection,
inversion, and left and right multiplication.  The Wilson action splits into a
positive-time part, its reflected part, and plaquette factors crossing the
reflection plane.  A crossing plaquette factor is a continuous central
function of a group element \(V\):
\[
  K_\beta(V)
  =
  \exp\left\{\frac{\beta}{N}
  \operatorname{Re}\operatorname{tr}V\right\}.
\]
Peter--Weyl theory expands it as
\begin{equation}
  K_\beta(V)
  =
  \sum_{\lambda\in\widehat G}c_\lambda(\beta)\chi_\lambda(V),
  \qquad
  c_\lambda(\beta)
  =
  \int_G K_\beta(V)\overline{\chi_\lambda(V)}\,\dd V .
  \label{eq:lattice-peter-weyl-crossing-kernel}
\end{equation}
For the Wilson plaquette kernel the coefficients entering the reflection
plane decomposition are nonnegative.  Matrix-coefficient orthogonality rewrites
each crossing contribution as a finite or absolutely convergent sum of terms
\[
  c_\lambda(\beta)\,
  F_{\lambda,i}[U_+]\,\overline{F_{\lambda,i}[\vartheta U_+]},
  \qquad c_\lambda(\beta)\geq0,
\]
after the links on the reflection plane have been integrated.  Multiplication
by the positive-time Boltzmann weight and summation over \(\lambda,i\) turns
\eqref{eq:lattice-reflection-positivity} into a sum of squared absolute values
with nonnegative coefficients.
\end{proof}

Reflection positivity produces a Hilbert space at finite lattice spacing by
quotienting \(\mathcal B_+\) by the null vectors of
\eqref{eq:lattice-reflection-positivity} and completing.  Translation by one
lattice unit in Euclidean time gives a positive transfer matrix
\[
  T_a=\ee^{-aH_a}.
\]
Thus the finite lattice theory has a genuine Hilbert-space interpretation with
a nonnegative lattice Hamiltonian \(H_a\).  The continuum question is whether a
family of such OS-positive finite-cutoff theories has a scaling limit whose
Schwinger functions satisfy the continuum hypotheses of
Chapter~\ref{ch:os-reconstruction}.

\section{Wilson Loops and the Strong-Coupling Area Law}

Let \(C\) be a closed oriented lattice path.  If \(C\) traverses the oriented
link \(\ell\) with sign \(\epsilon_\ell=\pm1\), set
\[
  U(C)=\prod_{\ell\in C}U_\ell^{\epsilon_\ell},
\]
with the product ordered along the path.  For an irreducible representation
\(R\) of \(G\) with character \(\chi_R\) and dimension \(d_R\), define the
normalized Wilson loop
\begin{equation}
  W_R(C)
  =
  d_R^{-1}\chi_R(U(C)).
  \label{eq:lattice-wilson-loop-definition}
\end{equation}
Gauge factors at adjacent endpoints cancel around the closed curve, so
\(W_R(C)\) is gauge invariant.

The strong-coupling expansion is the character expansion of each plaquette
Boltzmann factor.  Write
\begin{equation}
  \exp\left\{\frac{\beta}{N}
  \operatorname{Re}\operatorname{tr}U_P\right\}
  =
  c_0(\beta)
  \left[
    1+\sum_{\lambda\neq 0}d_\lambda u_\lambda(\beta)\chi_\lambda(U_P)
  \right],
  \qquad
  u_\lambda(\beta)=O(\beta^{n_\lambda}),
  \label{eq:lattice-character-expansion}
\end{equation}
where \(n_\lambda>0\) for every nontrivial representation.  The Haar
orthogonality relation
\begin{equation}
  \int_G
  R^\lambda(U)_{ij}R^\mu(U^{-1})_{kl}\,\dd U
  =
  \frac{\delta_{\lambda\mu}}{d_\lambda}
  \delta_{il}\delta_{jk}
  \label{eq:lattice-haar-orthogonality}
\end{equation}
implies a local constraint: after integrating a link, the representation
matrices carried by the plaquettes and by the Wilson loop must combine to the
trivial representation on that link.

For a fundamental Wilson loop in \(SU(N)\), the lowest nonzero contribution is
therefore obtained by filling a surface \(\Sigma\) with boundary \(C\) by
fundamental plaquettes.  If \(A_{\min}(C)\) is the minimal number of plaquettes
in such a surface, then
\begin{equation}
  \langle W_{\square}(C)\rangle_{\Lambda,\beta}
  =
  K_C\,
  u_{\square}(\beta)^{A_{\min}(C)}
  \bigl(1+O(\beta)\bigr)
  \label{eq:lattice-strong-coupling-leading-area}
\end{equation}
for sufficiently small \(\beta\), with \(K_C>0\) determined by boundary group
integrals and by the normalization of \(W_{\square}\).  More precisely, the
polymer expansion obtained from \eqref{eq:lattice-character-expansion}
converges absolutely when \(\beta\) is small enough; its polymer activities
obey a bound of the form
\[
  \sum_{\Gamma\ni P}|z(\Gamma)|\ee^{\alpha|\Gamma|}
  \leq \alpha
\]
for a positive \(\alpha\), after decreasing \(\beta\) if necessary.  This
Kotecky--Preiss type estimate gives constants \(C(\beta)\) and
\(\sigma(\beta)>0\) such that
\begin{equation}
  |\langle W_{\square}(C)\rangle_{\Lambda,\beta}|
  \leq
  C(\beta)\,
  \exp\{-\sigma(\beta)A_{\min}(C)\}
  \label{eq:lattice-strong-coupling-area-bound}
\end{equation}
uniformly in the finite volume, away from boundary effects.

For a rectangular loop with spatial length \(R=na\) and Euclidean time length
\(T=ma\), \eqref{eq:lattice-strong-coupling-area-bound} gives
\[
  \langle W_{\square}(C_{R,T})\rangle
  \sim
  \exp\{-\sigma_{\rm lat}(\beta)nm\}
  =
  \exp\{-[\sigma_{\rm lat}(\beta)R/a]\,T/a\}.
  \]
The transfer-matrix interpretation therefore gives a static-source potential
which grows linearly with separation in the strong-coupling lattice theory.
This is a rigorous confinement diagnostic at finite lattice cutoff for small
\(\beta\).  The four-dimensional continuum confinement problem asks for
control along a scaling trajectory with \(\beta\to\infty\).

\section{Scaling Trajectories and Dimensional Transmutation}

The continuum limit is not obtained by fixing \(\beta\) and sending
\(a\downarrow0\).  One must choose a trajectory \(g_0(a)\), equivalently
\(\beta(a)=N/g_0(a)^2\) in the trace normalization of this chapter, so that
dimensionless lattice observables scale in
a way compatible with finite physical observables.  For example, if the
continuum theory has a string tension \(\sigma_{\rm phys}\), a scaling
trajectory would satisfy
\[
  a^2\sigma_{\rm lat}(\beta(a))
  \longrightarrow
  a^2\sigma_{\rm phys}
  \quad\text{with}\quad
  a^2\sigma_{\rm phys}\downarrow0.
\]
Equivalently, the correlation length in lattice units must diverge.

Perturbation theory around the Gaussian ultraviolet fixed point predicts the
large-\(\beta\) relation between \(a\) and the lattice coupling.  For pure
\(SU(N)\), define \(\beta_0,\beta_1\) by
\begin{equation}
  \mu\frac{\dd g}{\dd\mu}
  =
  -\beta_0 g^3-\beta_1 g^5+O(g^7),
  \qquad
  \beta_0=\frac{11N}{48\pi^2},
  \qquad
  \beta_1=\frac{34N^2}{3(16\pi^2)^2}.
  \label{eq:lattice-pure-ym-beta-coefficients}
\end{equation}
Then the associated lattice scale \(\Lambda_L\) is characterized by
\begin{equation}
  a\Lambda_L
  =
  \exp\left\{-\frac{1}{2\beta_0g_0^2}\right\}
  (\beta_0g_0^2)^{-\beta_1/(2\beta_0^2)}
  \bigl(1+O(g_0^2)\bigr),
  \qquad
  g_0^2=\frac{2N}{\beta}.
  \label{eq:lattice-asymptotic-scaling}
\end{equation}
The numerical value of \(\Lambda_L\) depends on the lattice action and on the
definition of the bare coupling.  Ratios between \(\Lambda_L\) and scales in
other schemes, such as \(\Lambda_{\overline{\mathrm{MS}}}\), are determined by
matching a short-distance observable in the overlap of the two regularized
perturbation theories.

The formula \eqref{eq:lattice-asymptotic-scaling} is not a substitute for the
nonperturbative continuum construction.  It gives the perturbative tangent to
the scaling trajectory near \(g_0=0\).  A construction of four-dimensional
Yang--Mills theory must also show that the gauge-invariant Schwinger functions
or an equivalent algebraic object converge along such a trajectory and that
the reconstructed Hilbert-space theory has the desired spectral properties.

\section{Local Operators, Counterterms, and Lattice Symmetries}

The Wilson lattice action has exact gauge invariance, lattice translations,
hypercubic rotations and reflections, and charge conjugation when the
\(\theta\)-angle is set to zero.  The continuum local terms that can be
generated by the lattice regulator must respect these symmetries.  This gives
the Symanzik form of a local long-distance effective action:
\begin{equation}
  S_{\rm Sym}
  =
  \int \dd^4x\,
  \left[
    \mathcal E_0(a)
    +
    \frac{1}{4g_R(a)^2}
    \operatorname{tr}(F_{\mu\nu}F_{\mu\nu})
    +
    \sum_i a^{d_i-4}c_i(a)\mathcal O_i(x)
  \right]
  +
  \ii\theta\,\frac{1}{8\pi^2}\int\operatorname{tr}(F\wedge F).
  \label{eq:lattice-symanzik-effective-action}
\end{equation}
Here \(\mathcal E_0(a)\) is a vacuum-energy density, \(d_i>4\), and
\(\mathcal O_i\) are gauge-invariant local operators compatible with the
lattice symmetries.  In pure \(SU(N)\) Yang--Mills, with CP imposed, the only
nontrivial gauge-invariant dimension-four term is
\(\operatorname{tr}(F_{\mu\nu}F_{\mu\nu})\).  If CP is not imposed, the
topological density \(\operatorname{tr}(F\wedge F)\) supplies the separate
\(\theta\)-parameter.  Gauge invariance forbids a gluon mass term, and
hypercubic symmetry makes the coefficient of \(F_{\mu\nu}F_{\mu\nu}\) the same
for every coordinate two-plane on an isotropic lattice.

\begin{definition}[Tuning for a lattice continuum limit]
A lattice parameter is tuned when it must be chosen as a function of \(a\) in
order to keep a specified renormalized observable finite and nonzero as
\(a\downarrow0\).  Parameters multiplying operators with dimension below four
are relevant at the Gaussian fixed point; parameters multiplying
dimension-four operators are marginal at tree level and must be classified by
the beta function and by symmetry.
\end{definition}

For pure CP-even \(SU(N)\) Yang--Mills on an isotropic Wilson lattice, this
power-counting and symmetry analysis leaves one running parameter, \(g_0(a)\).
The mass gap, glueball spectrum, and string tension are then outputs of the
continuum theory, if the continuum limit exists.  Different gauge-invariant
lattice actions with the same symmetries change the coefficients of irrelevant
operators in \eqref{eq:lattice-symanzik-effective-action} and change the
relation between their lattice scales.  Symanzik improvement chooses additional
irrelevant plaquette and loop terms so that the leading lattice artifacts begin
at higher powers of \(a\).

For theories whose ultraviolet flow is not asymptotically free, the preceding
one-parameter mechanism may fail to produce an interacting continuum theory.
Four-dimensional scalar \(\phi^4\) theory and four-dimensional QED are the
basic cautionary examples: perturbation theory points toward a Landau-pole
obstruction, and rigorous triviality results constrain large classes of
lattice scalar limits.  Those results do not settle every possible
UV-complete theory that has the same formal perturbative expansion around the
Gaussian point.  The monograph therefore treats triviality and asymptotic
non-freedom as theory-specific continuum-limit questions, not as consequences
of formal power series alone.

\section[Theta-Term and Lattice Topology]{The \(\theta\)-Term and Lattice Topology}

For a smooth continuum \(SU(N)\) connection on a closed four-manifold,
\[
  Q[A]
  =
  \frac{1}{8\pi^2}\int \operatorname{tr}(F\wedge F)
  \in\mathbb Z
\]
is the second Chern number.  The Euclidean action may include
\(\ii\theta Q[A]\), and the integer-valuedness of \(Q\) gives
\(2\pi\)-periodicity in \(\theta\).

At finite lattice spacing, a generic lattice gauge field is only an assignment
of group elements to links.  There is no literal smooth bundle connection from
which \(Q\) can be read without an additional construction.  A plaquette
discretization of the density \(\operatorname{tr}(F\wedge F)\) has the correct
classical limit on smooth fields, but it is not integer-valued on arbitrary
lattice fields.  Consequently \(\exp(\ii\theta Q_{\rm plaq})\) need not be
\(2\pi\)-periodic at finite \(a\).

There are several regulator-level ways to define topological charge more
carefully.  A geometric definition interpolates admissible lattice fields to a
bundle over the torus and evaluates the corresponding Chern number.  With
Ginsparg--Wilson fermions, the index of the lattice Dirac operator gives an
integer
\[
  Q_{\rm GW}
  =
  \operatorname{Tr}\,\gamma_5\left(1-\frac{aD_{\rm GW}}{2}\right)
\]
under the usual admissibility hypotheses.  Gradient-flow definitions smooth
the gauge field before evaluating a lattice density and are designed to have a
controlled continuum limit.  These choices agree in the continuum scaling
regime when the required estimates hold, but they are distinct finite-cutoff
definitions.  Thus the \(\theta\)-angle is an additional parameter whose
nonperturbative lattice definition requires its own regulator convention.

\section{Fermions: Doubling and Chiral Lattice Actions}

A vectorlike gauge theory with fermions requires a lattice Dirac operator
coupled to the link variables.  The naive free lattice Dirac operator is
\begin{equation}
  D_{\rm n}(p)
  =
  \frac{\ii}{a}\sum_{\mu=1}^4\gamma_\mu\sin(ap_\mu),
  \qquad
  p_\mu\in[-\pi/a,\pi/a).
  \label{eq:lattice-naive-dirac-symbol}
\end{equation}
Its zeros occur whenever every \(ap_\mu\) is \(0\) or \(\pi\).  In four
dimensions this gives \(2^4=16\) light species near the corners of the
Brillouin zone.

\begin{figure}[H]
  \centering
  \resizebox{0.72\textwidth}{!}{%
  \begin{tikzpicture}[x=1cm,y=1cm,>=Stealth,every node/.style={font=\scriptsize}]
    \draw[->] (-3.2,0)--(3.35,0) node[right] {\(ap_1\)};
    \draw[->] (0,-3.2)--(0,3.35) node[above] {\(ap_2\)};
    \draw (-3.1416,-0.08)--(-3.1416,0.08) node[below=4pt] {\(-\pi\)};
    \draw (3.1416,-0.08)--(3.1416,0.08) node[below=4pt] {\(\pi\)};
    \draw (-0.08,-3.1416)--(0.08,-3.1416) node[left=4pt] {\(-\pi\)};
    \draw (-0.08,3.1416)--(0.08,3.1416) node[left=4pt] {\(\pi\)};
    \draw[blue!55!black,thick] (-3.1416,-3.1416) rectangle (3.1416,3.1416);
    \foreach \x in {-3.1416,0,3.1416} {
      \foreach \y in {-3.1416,0,3.1416} {
        \fill[red!75!black] (\x,\y) circle (2.7pt);
      }
    }
    \node[align=center] at (0,-3.85)
      {zeros of \(\sin(ap_\mu)\) in two dimensions;\\
       four dimensions have \(2^4\) corners};
  \end{tikzpicture}
  }
  \caption{The naive lattice derivative vanishes at every Brillouin-zone
  corner.  Each corner supplies a continuum Dirac cone, producing species
  doubling.}
  \label{fig:lattice-fermion-doubling}
\end{figure}

\begin{theorem}[Nielsen--Ninomiya obstruction, structural form]
Let \(D(p)\) be a free lattice Dirac operator on a periodic hypercubic lattice
which is local in position space, translationally invariant, has the correct
single-particle continuum behavior near \(p=0\), anticommutes with a fixed
\(\gamma_5\), and satisfies the usual Hermiticity condition giving a real
spectrum for \(\gamma_5D\).  Then the zeros of \(D(p)\) in the Brillouin torus
come with total chirality zero.  In particular, these hypotheses do not allow
a single Weyl fermion.
\end{theorem}

\begin{proof}
The normalized symbol of \(D(p)\) defines a map from the Brillouin torus with
small balls around its zeros removed to a sphere of Dirac matrices.  The
chirality of a zero is the local winding number of this map around that zero.
The torus has no boundary.  Applying degree additivity to the boundary of the
punctured torus gives that the sum of all local winding numbers is zero.  A
single continuum Weyl fermion would contribute a nonzero total winding number,
contradicting this identity.
\end{proof}

Standard lattice fermion constructions evade different hypotheses of this
theorem.
\begin{description}[leftmargin=2.9em,style=nextline]
\item[Wilson fermions]
Add \(r a\,\Delta/2\), giving doublers masses of order \(1/a\).  The Wilson
term breaks chiral symmetry and produces additive mass renormalization, so the
chiral limit requires tuning \(m_0(a)\) to a critical value \(m_c(a)\).
\item[Staggered fermions]
Diagonalize part of the spin structure and reduce the number of tastes.  The
residual taste symmetry and rooting questions must be treated as finite-cutoff
regularization issues.
\item[Domain-wall fermions]
Place chiral modes on separated boundaries of an auxiliary fifth direction.
Finite fifth-dimensional size produces residual chiral-symmetry breaking.
\item[Overlap and Ginsparg--Wilson fermions]
Replace \(\{\gamma_5,D\}=0\) by
\(\gamma_5D+D\gamma_5=aD\gamma_5D\), giving an exact modified lattice chiral
symmetry and an index theorem at finite cutoff under admissibility hypotheses.
\end{description}
For vectorlike QCD, these formulations provide regulator choices for
gauge-invariant correlation functions.  Their continuum equivalence requires
the same kind of scaling-limit control as in the pure gauge theory, together
with the relevant mass and chiral-symmetry tunings.

\section{Status of the Four-Dimensional Continuum Construction}

\begin{hypothesis}[Four-dimensional lattice Yang--Mills continuum limit]
\label{hyp:lattice-yang-mills-continuum-limit}
Let \(G=SU(N)\) and let the Wilson lattice measures
\(\dd\mu_{\Lambda,\beta}\) be defined by
\eqref{eq:lattice-partition-function}.  There exists a scaling trajectory
\(\beta(a)\to\infty\) and an infinite-volume limit such that the
gauge-invariant Schwinger functions of local and loop observables converge,
after the required field and operator normalizations, to Schwinger functions
satisfying the continuum Osterwalder--Schrader hypotheses.  The reconstructed
Hilbert-space theory has a unique vacuum and a positive mass gap.
\end{hypothesis}

This hypothesis is the constructive content of four-dimensional pure
Yang--Mills theory.  The finite lattice theory is rigorous; the hypothesis is
the assertion that the rigorous finite-cutoff objects have the required
continuum limit and spectral gap.  Renormalization-group analyses of lattice
gauge theory, including Ba{\l}aban's multiscale program, prove substantial
estimates toward such a construction.  Other constructive results in related
models and dimensions give additional control of particular mechanisms.  A
complete proof of Hypothesis~\ref{hyp:lattice-yang-mills-continuum-limit} for
four-dimensional nonabelian Yang--Mills theory remains the Yang--Mills
existence and mass-gap problem.

The strong-coupling area law is a theorem inside the small-\(\beta\) lattice
regime.  The weak-coupling continuum regime lies at \(\beta\to\infty\).  For
four-dimensional pure \(SU(2)\) and \(SU(3)\) Wilson gauge theory, numerical
and analytic evidence support a confining continuum limit connected to the
strong-coupling phase without a physical deconfining transition at zero
temperature.  For other groups and lattice actions, bulk lattice-artifact
transitions may occur and must be distinguished from continuum physics.  The
logical role of the lattice construction is therefore precise: it supplies a
finite, gauge-invariant nonperturbative regulator and a clear continuum-limit
problem; it does not by itself solve the mass-gap problem.

Covariant gauge fixing, Faddeev--Popov ghosts, BRST cohomology, and the BV
formalism developed in the following chapters are indispensable perturbative
and cohomological tools.  Their nonperturbative identification with the
continuum theory constructed from the lattice rests on two further inputs:
first, the lattice continuum limit in
Hypothesis~\ref{hyp:lattice-yang-mills-continuum-limit}; second, agreement of
short-distance gauge-invariant observables with the covariant perturbative
expansion order by order after matching the renormalized coupling and local
operator normalizations.
